% GENERATED FILE. Edit canonical lessons, then run npm run generate:books.
% canonical-source-hash: fnv1a64:1db20f88085a878f
% canonical-lessons: RU-C113-game, RU-C113-film, RU-C113-song, RU-C113-drawing, RU-C113-life

\chapter{Game, Film, Song, Drawing, Life}
\label{ch:ru-game-film-song-drawing-life}
\hlchaptermodality{113}

\begin{chapteropening}
\textbf{By the end of this chapter:} \emph{I can say a game, a film, a song, a drawing and life.}

Complete the last lesson of chapter 113: I can say a game, a film, a song, a drawing and life.
\end{chapteropening}

\section[\texorpdfstring{\ru{игра} (igrá)}{игра (igrá)}]{\ru{игра} (igrá) — a game}

\label{lesson:RU-C113-game}



\begin{quote}
Before the new one: say the Russian for music, then the Russian for sport.
\end{quote}

\subsection*{You'll want to know: \ru{игра}}
\textbf{\ru{игра}} (\emph{igrá}) — \textquotedblleft{}a game\textquotedblright{}.

\begin{cousinweb}
\textbf{\ru{игра}} $\leftarrow$ Proto-Slavic \emph{*jĭgra} — inherited. Vasmer relates it to Lithuanian \emph{áikštytis}, \textquotedblleft{}to be playful\textquotedblright{}, and Sanskrit \emph{ejati}, \textquotedblleft{}stirs, moves\textquotedblright{}.
\end{cousinweb}

\begin{grammarlens}[title={What you've built}]
One more word for this chapter.
\end{grammarlens}

\subsection*{Your turn}
\begin{itemize}
\raggedright
  \item \emph{Say it:} \emph{igrá}
  \item \emph{Say it:} \emph{igrá}, once more
  \item \emph{Say it:} say \emph{igrá}
  \item \emph{Recall:} say the Russian for music, then the Russian for sport, then say \emph{igrá} again
\end{itemize}

\begin{tcolorbox}[breakable,colback=teal!4,colframe=teal!35!black,title={Before you move on}]
What does \ru{игра} mean? (\textquotedblleft{}A game\textquotedblright{}.) Say it once more.
\end{tcolorbox}
\section[\texorpdfstring{\ru{фильм} (film)}{фильм (film)}]{\ru{фильм} (film) — a film}

\label{lesson:RU-C113-film}



\begin{quote}
Before the new one: say the Russian for sport, then the Russian for a game.
\end{quote}

\subsection*{You'll want to know: \ru{фильм}}
\textbf{\ru{фильм}} (\emph{film}) — \textquotedblleft{}a film\textquotedblright{}.

\begin{cousinweb}
\textbf{\ru{фильм}}, older \emph{\ru{фильма}}, comes from English \emph{film}, \textquotedblleft{}a thin skin, a film\textquotedblright{} (Vasmer).
\end{cousinweb}

\begin{grammarlens}[title={What you've built}]
One more word for this chapter.
\end{grammarlens}

\subsection*{Your turn}
\begin{itemize}
\raggedright
  \item \emph{Say it:} \emph{film}
  \item \emph{Say it:} \emph{film}, once more
  \item \emph{Say it:} say \emph{film}
  \item \emph{Recall:} say the Russian for sport, then the Russian for a game, then say \emph{film} again
\end{itemize}

\begin{tcolorbox}[breakable,colback=teal!4,colframe=teal!35!black,title={Before you move on}]
What does \ru{фильм} mean? (\textquotedblleft{}A film\textquotedblright{}.) Say it once more.
\end{tcolorbox}
\section[\texorpdfstring{\ru{песня} (pésnya)}{песня (pésnya)}]{\ru{песня} (pésnya) — a song}

\label{lesson:RU-C113-song}



\begin{quote}
Before the new one: say the Russian for a game, then the Russian for a film.
\end{quote}

\subsection*{You'll want to know: \ru{песня}}
\textbf{\ru{песня}} (\emph{pésnya}) — \textquotedblleft{}a song\textquotedblright{}.

\begin{cousinweb}
\textbf{\ru{песня}} comes from \emph{\ru{петь}}, \textquotedblleft{}to sing\textquotedblright{}, \emph{\ru{пою}}, \textquotedblleft{}I sing\textquotedblright{} (Vasmer): Old Russian \emph{\ru{песнь}}, \textquotedblleft{}what is sung\textquotedblright{}.
\end{cousinweb}

\begin{grammarlens}[title={What you've built}]
One more word for this chapter.
\end{grammarlens}

\subsection*{Your turn}
\begin{itemize}
\raggedright
  \item \emph{Say it:} \emph{pésnya}
  \item \emph{Say it:} \emph{pésnya}, once more
  \item \emph{Say it:} say \emph{pésnya}
  \item \emph{Recall:} say the Russian for a game, then the Russian for a film, then say \emph{pésnya} again
\end{itemize}

\begin{tcolorbox}[breakable,colback=teal!4,colframe=teal!35!black,title={Before you move on}]
What does \ru{песня} mean? (\textquotedblleft{}A song\textquotedblright{}.) Say it once more.
\end{tcolorbox}
\section[\texorpdfstring{\ru{рисунок} (risúnok)}{рисунок (risúnok)}]{\ru{рисунок} (risúnok) — a drawing}

\label{lesson:RU-C113-drawing}



\begin{quote}
Before the new one: say the Russian for a film, then the Russian for a song.
\end{quote}

\subsection*{You'll want to know: \ru{рисунок}}
\textbf{\ru{рисунок}} (\emph{risúnok}) — \textquotedblleft{}a drawing\textquotedblright{}.

\begin{cousinweb}
\textbf{\ru{рисунок}} — recorded from 1705; it comes from Polish \emph{rysunek}, \textquotedblleft{}a drawing\textquotedblright{}, from \emph{rysować}, \textquotedblleft{}to draw\textquotedblright{}, with an ending that reflects German \emph{-ung} (Vasmer). Polish \emph{rysować} goes back to German \emph{reißen}, \textquotedblleft{}to draw, to sketch\textquotedblright{}.
\end{cousinweb}

\begin{grammarlens}[title={What you've built}]
One more word for this chapter.
\end{grammarlens}

\subsection*{Your turn}
\begin{itemize}
\raggedright
  \item \emph{Say it:} \emph{risúnok}
  \item \emph{Say it:} \emph{risúnok}, once more
  \item \emph{Say it:} say \emph{risúnok}
  \item \emph{Recall:} say the Russian for a film, then the Russian for a song, then say \emph{risúnok} again
\end{itemize}

\begin{tcolorbox}[breakable,colback=teal!4,colframe=teal!35!black,title={Before you move on}]
What does \ru{рисунок} mean? (\textquotedblleft{}A drawing\textquotedblright{}.) Say it once more.
\end{tcolorbox}
\section[\texorpdfstring{\ru{жизнь} (zhizn)}{жизнь (zhizn)}]{\ru{жизнь} (zhizn) — life}

\label{lesson:RU-C113-life}



\begin{quote}
Before the new one: say the Russian for a song, then the Russian for a drawing.
\end{quote}

\subsection*{You'll want to know: \ru{жизнь}}
\textbf{\ru{жизнь}} (\emph{zhizn}) — \textquotedblleft{}life\textquotedblright{}.

\begin{cousinweb}
\textbf{\ru{жизнь}} comes from \emph{\ru{жить}}, \textquotedblleft{}to live\textquotedblright{} (Vasmer), with the old noun suffix \emph{-\ru{знь}}, as in \emph{\ru{болезнь}}, \textquotedblleft{}illness\textquotedblright{}.
\end{cousinweb}

\begin{grammarlens}[title={What you've built}]
One more word for this chapter.
\end{grammarlens}

\subsection*{Your turn}
\begin{itemize}
\raggedright
  \item \emph{Say it:} \emph{zhizn}
  \item \emph{Say it:} \emph{zhizn}, once more
  \item \emph{Say it:} say \emph{zhizn}
  \item \emph{Recall:} say the Russian for a song, then the Russian for a drawing, then say \emph{zhizn} again
\end{itemize}

\begin{tcolorbox}[breakable,colback=teal!4,colframe=teal!35!black,title={Before you move on}]
What does \ru{жизнь} mean? (\textquotedblleft{}Life\textquotedblright{}.) Say it once more.
\end{tcolorbox}
